%% IEEE Conference Template — IEEEtran class
%% Compile: pdflatex -> bibtex -> pdflatex -> pdflatex
\documentclass[conference]{IEEEtran}
\IEEEoverridecommandlockouts

% ── Packages ──────────────────────────────────────────────────────────────────
\usepackage{cite}
\usepackage{amsmath,amssymb,amsfonts}
\usepackage{algorithmic}
\usepackage{graphicx}
\usepackage{textcomp}
\usepackage{xcolor}
\usepackage{booktabs}
\usepackage{multirow}
\usepackage{url}
\usepackage{microtype}

\def\BibTeX{{\rm B\kern-.05em{\sc i\kern-.025em b}\kern-.08em
    T\kern-.1667em\lower.7ex\hbox{E}\kern-.125emX}}

% ── Document ──────────────────────────────────────────────────────────────────
\begin{document}

% ── Title ─────────────────────────────────────────────────────────────────────
\title{Training-Set Diversity Does Not Eliminate Ancestry-Associated\\
Reference Bias in DNA Foundation Models}

% ── Authors ───────────────────────────────────────────────────────────────────
\author{
\IEEEauthorblockN{Thackshanaramana B}
\IEEEauthorblockA{\textit{Department of Computational Intelligence} \\
\textit{College of Engineering and Technology} \\
SRM Institute of Science and Technology \\
Kattankulathur, Tamil Nadu, India \\
thackshanaramana0@gmail.com}
}

\maketitle

% ── Abstract ──────────────────────────────────────────────────────────────────
\begin{abstract}
DNA foundation models are increasingly used for genomic variant interpretation,
yet whether diversity in pretraining data is sufficient to mitigate
ancestry-associated scoring bias remains unclear. Here, we examine reference
allele preference in Nucleotide Transformer using ancestry-enriched common
variants from gnomAD and compare models pretrained on the human reference genome
and on 3{,}202 genetically diverse human genomes. Both models exhibit a
significant ancestry-associated scoring asymmetry, with European-enriched
variants receiving greater reference allele preference than African-enriched
variants. Strikingly, the EUR--AFR bias gap changes by only 1.2\% following
diverse-genome pretraining. Neither population allele frequency nor local
genomic variant density explains the observed difference, indicating that the
effect is not simply driven by variant prevalence or regional sequence variation.
These results show that substantially increasing genomic diversity during
pretraining does not eliminate ancestry-associated reference bias in Nucleotide
Transformer and point to a persistent reference-centered component in DNA
foundation-model scoring. More representative genomic representations, including
pangenome-aware approaches, may therefore be necessary to address
ancestry-associated asymmetries in genomic foundation models.
\end{abstract}

\begin{IEEEkeywords}
DNA foundation models, reference allele bias, ancestry bias, health equity,
Nucleotide Transformer, masked language model, variant effect prediction,
pangenome, gnomAD, population genomics
\end{IEEEkeywords}

% ── I. Introduction ───────────────────────────────────────────────────────────
\section{Introduction}

The human reference genome GRCh38 (hg38) is the universal coordinate system of
clinical and computational genomics. However, hg38 is not population-neutral:
approximately 70\% of its sequence derives from a single donor (RP11), whose
ancestry has been characterised as African-European admixed \cite{t2t}.
Crucially, hg38 is a \emph{single linear haplotype}---it encodes one path
through human sequence diversity, anchoring all downstream analyses to one
set of reference alleles regardless of the donor's precise ancestry. As of 2023, 86.5\% of participants in the GWAS Catalog were
of European descent \cite{diversity_gap}, and non-European patients experience
a ${\sim}23\%$ increase in variants of uncertain significance (VUS) compared
to Europeans \cite{vus_burden}---a disparity partly attributed to reference
genome bias propagating through analytical tools.

DNA sequence foundation models represent the newest class of such tools. These
large transformer architectures---pretrained on genomic sequences via
self-supervised objectives---have become central to computational genomics,
with applications spanning regulatory element prediction, variant effect scoring,
and chromatin accessibility modelling \cite{dalla_torre_nt, dart_eval, gpn_benegas}.
The Nucleotide Transformer (NT) \cite{dalla_torre_nt}, published in
\emph{Nature Methods} (2025), exemplifies this class with two 500M-parameter
models: NT-hg38, pretrained on the single human reference genome, and NT-1000G,
pretrained on 3{,}202 genetically diverse genomes from the 1000 Genomes Project
(1000G) spanning 26 populations including ${\sim}660$ African samples (26\%).
The NT-1000G model was explicitly designed to \emph{``encode a better
representation of human genetic variation''} \cite{dalla_torre_nt}.

Reference allele bias is a specific and measurable failure mode for these
models. Under MLM scoring, the model receives a masked k-mer context and
predicts the probability of each vocabulary token at the masked position. The
\emph{bias score} is defined as:
\begin{equation}
  S = \log P(\text{ref-kmer} \mid \text{context})
      - \log P(\text{alt-kmer} \mid \text{context})
  \label{eq:bias}
\end{equation}
Positive $S$ means the model assigns higher probability to the reference
allele; negative $S$ means it prefers the alternate. A population-fair model
should score variants with $S \approx 0$ when the alternate allele is common
in the sampled population, regardless of ancestry. Deviation from this
expectation quantifies the model's reference allele preference.

Pathak et al. \cite{pathak_2024} demonstrated pervasive ancestry bias in 52
classical variant effect predictors (VEPs) and in protein language models,
finding that African-ancestry individuals are disproportionately assigned high
pathogenicity scores for benign variants. However, no study has examined whether
\emph{DNA sequence foundation models} exhibit analogous ancestry-asymmetric
reference bias, nor whether diverse training data (NT-1000G's design principle)
corrects it.

We address three specific questions:
\begin{enumerate}
  \item Do NT models show significant positive reference allele bias for
        African-enriched common variants?
  \item Is this bias symmetric across ancestries when within-population allele
        frequency is matched?
  \item Does diverse training (NT-1000G) correct any observed asymmetry?
\end{enumerate}

% ── II. Related Work ──────────────────────────────────────────────────────────
\section{Related Work}

\subsection{Ancestry Bias in Variant Effect Prediction}
Pathak et al.~\cite{pathak_2024} examined 52 VEPs across 14 ancestry groups
for missense variants, finding that tools relying on population or clinical
training data assign disproportionately high pathogenicity scores to variants
enriched in underrepresented populations. Protein language models (EVE, ESM-1v)
show reduced but non-zero bias. Our work is the first to apply this framework
to DNA sequence foundation models.

\subsection{Nucleotide Transformer}
Dalla-Torre et al.~\cite{dalla_torre_nt} introduced NT-hg38 and NT-1000G,
evaluating downstream task performance on genomic benchmarks. The NT-1000G model
is trained on 3{,}202 phased genomes from 27 geographically diverse populations
totalling 19.2 trillion nucleotides. The authors report improved performance on
some tasks; no ancestry-stratified analysis of the internal bias scoring function
is performed.

\subsection{DNA LM Benchmarks}
DART-Eval \cite{dart_eval} (NeurIPS 2024) benchmarks five DNA language models on
regulatory tasks including African caQTL datasets, measuring task accuracy rather
than internal scoring bias. Feng et al.~\cite{benchmark_dna_lm} benchmark
five DNA LMs on genomic classification and variant tasks; population stratification
of bias is not examined. Neither work measures reference allele preference as a
function of variant ancestry.

\subsection{MLM Bias Scoring for Variant Effect}
Benegas et al.~\cite{gpn_benegas} applied MLM bias scoring genome-wide for
variant effect prediction via the Genomic Pre-trained Network (GPN), demonstrating
that DNA language models can serve as unsupervised VEPs. No ancestry-stratified
analysis of scoring bias is performed.

\subsection{Reference Genome Diversity}
The T2T-CHM13 assembly and the Human Pangenome Reference Consortium
(HPRC) \cite{pangenome} have produced ancestry-spanning reference sequences.
No DNA foundation model trained on pangenome references has been publicly
released, leaving open the question of whether pangenome training would reduce
ancestry-specific scoring bias.

\noindent\textbf{Gap:} No prior work has measured reference allele bias in DNA
foundation models stratified by ancestry-enriched variant sets, nor tested
whether diverse training corrects it.

% ── III. Methods ──────────────────────────────────────────────────────────────
\section{Methods}

\subsection{Variant Selection from gnomAD v4}

We queried gnomAD v4~\cite{gnomad_v4} via the public GraphQL API for all
chromosome 22 SNVs with population-stratified allele frequencies. Two variant
sets were constructed using symmetric criteria:

\noindent\textbf{African-enriched (AFR):} $\text{AF}_\text{AFR} \geq 0.10$,
$\text{AF}_\text{EUR} \leq 0.02$, fold-enrichment
$(\text{AF}_\text{AFR}/\text{AF}_\text{EUR}) \geq 5\times$.
Yielded $N{=}7{,}684$ variants; within-population AF mean 0.179 (IQR 0.120--0.225).

\noindent\textbf{European-enriched (EUR):} $\text{AF}_\text{EUR} \geq 0.10$,
$\text{AF}_\text{AFR} \leq 0.02$.
Yielded $N{=}509$ variants; within-population AF mean 0.171 (IQR 0.118--0.196).

The two sets are matched for within-population allele frequency magnitude
(Wilcoxon $p{=}0.31$). They differ substantially in global minor allele
frequency: AFR mean 0.031, EUR mean 0.111. We chose symmetric thresholds to
ensure the comparison is driven by ancestry enrichment, not by frequency
magnitude.

\subsection{Models}

Both models were loaded from HuggingFace in CPU-only mode
(CUDA disabled via \texttt{CUDA\_VISIBLE\_DEVICES=""} before import,
required to avoid GPU architecture incompatibility with older
accelerators):
\begin{itemize}
  \item \textbf{NT-hg38:} \texttt{InstaDeepAI/nucleotide-transformer-500m-human-ref} ---
        trained on the single GRCh38 reference genome.
  \item \textbf{NT-1000G:} \texttt{InstaDeepAI/nucleotide-transformer-500m-1000g} ---
        trained on 3{,}202 genetically diverse human genomes from the 1000G
        Project (26\% African, 20\% European, 20\% East Asian, 20\% South Asian,
        14\% Admixed American).
\end{itemize}

\subsection{MLM Bias Scoring}

For each variant at chromosomal position $\text{pos}$, a 512-token window was
extracted from GRCh38 chromosome 22, centred on the variant. Sequences were
tokenised using the model's 6-mer non-overlapping vocabulary. The k-mer
containing the variant nucleotide was identified and its token masked. The model
computed log-softmax probabilities over the vocabulary at the masked position,
and the bias score $S$ was computed per Eq.~\eqref{eq:bias}. All scoring was
performed in FP32 without gradient computation. Total scoring calls: 16{,}386
(7{,}684 AFR + 509 EUR $\times$ 2 models); all completed with status \texttt{ok}.

\subsection{Statistical Analysis}

Between-group comparisons used Welch's $t$-test (unequal variance and sample
size). Effect sizes are reported as Cohen's $d$. Bootstrap 95\% confidence
intervals (10{,}000 iterations with \texttt{numpy.random.default\_rng(42)})
were computed for the EUR$-$AFR mean difference. Balanced resampling (1,000
iterations; 509 AFR drawn uniformly from the full pool vs.\ 509 EUR) was used
to confirm robustness to sample-size imbalance. Within-group allele frequency
calibration was assessed by Pearson correlation $r(\text{AF}, S)$ and Spearman
$\rho$; Fisher's $z$-test compared correlations across groups. Local variant
density was computed as the count of gnomAD common variants
($\text{AF}_\text{AFR}{\geq}0.05$ or $\text{AF}_\text{EUR}{\geq}0.05$) within
$\pm$500 bp of each focal variant.

% ── IV. Results ───────────────────────────────────────────────────────────────
\section{Results}

\subsection{Both Models Show Strong Positive Reference Bias for\\African-Enriched Variants}

Table~\ref{tab:descriptive} summarises bias score distributions. NT-hg38 scores
7{,}684 African-enriched variants with mean $+0.582$ (${\pm}2.227$;
$t$-test vs.\ zero: $p{=}2.4{\times}10^{-112}$; 59.7\% of variants positive).
NT-1000G yields mean $+0.508$ ($p{=}3.7{\times}10^{-86}$; 58.4\% positive).
These effects are large in terms of $p$-value but modest in absolute magnitude:
the model still systematically prefers the reference allele even when the
alternate allele is present at 10--50\% frequency in African populations.

\begin{table}[t]
\caption{Bias score descriptive statistics per group and model}
\label{tab:descriptive}
\centering
\renewcommand{\arraystretch}{1.15}
\begin{tabular}{@{}llrrrr@{}}
\toprule
Model & Group & $N$ & Mean & SD & $p$ (vs.\ 0) \\
\midrule
\multirow{2}{*}{NT-hg38}
  & AFR & 7,684 & $+$0.582 & 2.227 & $2.4{\times}10^{-112}$ \\
  & EUR &   509 & $+$1.010 & 2.280 & $1.5{\times}10^{-21}$ \\
\midrule
\multirow{2}{*}{NT-1000G}
  & AFR & 7,684 & $+$0.508 & 2.237 & $3.7{\times}10^{-86}$ \\
  & EUR &   509 & $+$0.941 & 2.298 & $7.0{\times}10^{-19}$ \\
\bottomrule
\end{tabular}
\end{table}

\subsection{Bias Is Ancestry-Asymmetric at Matched Within-Population Frequency}

Table~\ref{tab:between_group} shows that European-enriched variants receive
significantly higher positive bias than African-enriched variants across both
models. The effect replicates independently: $\Delta{=}{+}0.428$
($p{=}4.7{\times}10^{-5}$, $d{=}0.19$) for NT-hg38 and $\Delta{=}{+}0.433$
($p{=}4.3{\times}10^{-5}$, $d{=}0.19$) for NT-1000G. Bootstrap 95\% CI for
the NT-1000G gap is $[{+}0.231,{+}0.640]$, entirely above zero.

To rule out sample-size asymmetry as a confound (7,684 AFR vs.\ 509 EUR),
we performed balanced resampling: 1,000 iterations each drawing 509 variants
uniformly at random from the AFR pool and comparing to the 509-variant EUR set.
The median EUR$-$AFR gap across iterations is $+0.429$
(95\% CI $[{+}0.243,{+}0.597]$) for NT-hg38 and $+0.428$
($[{+}0.233,{+}0.623]$) for NT-1000G, with median Cohen's $d{=}0.190$ in both
cases. The gap exceeds zero in \textbf{all 1,000 iterations} and reaches
$p{<}0.05$ in 94.8\% (NT-hg38) and 94.3\% (NT-1000G) of iterations. The result
is not attributable to sample-size asymmetry.

Pooling all 8,193 scored variants and ranking by bias score (descending),
EUR-enriched variants occupy a mean score percentile of 55.8 vs.\ 49.6 for
AFR-enriched variants (NT-hg38; NT-1000G: 55.3 vs.\ 49.6). EUR variants
represent 8.2\% of the top decile despite comprising only 6.2\% of the pool
(1.32$\times$ enrichment), while AFR variants are proportionally displaced
toward lower-scoring positions. This cumulative displacement, operating across
thousands of common variants simultaneously, illustrates that even a modest
per-variant gap can produce a systematic population-level scoring advantage.

This direction is counterintuitive. European-enriched variants have 3.6$\times$
higher global minor allele frequency (0.111 vs.\ 0.031); a model tracking
global frequency would assign \emph{lower} positive bias to European-enriched
variants (their alternate allele is more globally common). The opposite is
observed. This rules out global allele frequency as the explanation.

\begin{table}[t]
\caption{Between-group comparison: EUR vs.\ AFR per model}
\label{tab:between_group}
\centering
\renewcommand{\arraystretch}{1.15}
\begin{tabular}{@{}lrrrr@{}}
\toprule
Model & $\Delta$ (EUR$-$AFR) & $t$ & $p$ & Cohen's $d$ \\
\midrule
NT-hg38  & $+$0.428 & $-$4.10 & $4.7{\times}10^{-5}$ & 0.19 \\
NT-1000G & $+$0.433 & $-$4.12 & $4.3{\times}10^{-5}$ & 0.19 \\
\bottomrule
\end{tabular}
\end{table}

\subsection{Training Diversity Leaves the Ancestry Gap Unchanged}

Table~\ref{tab:between_model} shows the effect of switching from NT-hg38 to
NT-1000G. For African-enriched variants, diversity training reduces mean bias
by 12.6\% in absolute terms (0.074 log-likelihood units; $p{=}0.041$,
$d{=}0.033$). For European-enriched variants, the reduction is 6.8\%
($p{=}0.635$, not significant).

The ancestry gap itself changes from $\Delta{=}0.428$ (NT-hg38) to
$\Delta{=}0.433$ (NT-1000G): a change of $\mathbf{+0.005}$, corresponding to a
\textbf{1.2\% change in the gap}. Adding 3{,}202 diverse training genomes---26\%
of which are African---modestly shifts absolute bias levels but leaves the
relative ancestry asymmetry statistically and practically unchanged.

\begin{table}[t]
\caption{Effect of diversity training (NT-hg38 $\rightarrow$ NT-1000G)}
\label{tab:between_model}
\centering
\renewcommand{\arraystretch}{1.15}
\begin{tabular}{@{}lrrrl@{}}
\toprule
Group & Reduction & $p$ & $d$ & Gap change \\
\midrule
AFR & 12.6\% & 0.041 & 0.033 & \multirow{2}{*}{$+$0.005 (1.2\%)} \\
EUR &  6.8\% & 0.635 & 0.030 & \\
\bottomrule
\end{tabular}
\end{table}

\subsection{Neither Model Is Calibrated to Population Allele Frequency}

Within the African-enriched group, the Pearson correlation between
$\text{AF}_\text{AFR}$ and bias score is $r{=}{-}0.053$ ($p{=}3.5{\times}10^{-6}$,
NT-hg38) and $r{=}{-}0.058$ ($p{=}4.2{\times}10^{-7}$, NT-1000G), explaining
$R^2{=}$0.28\% and 0.33\% of variance respectively. Within the European-enriched
group, $r{=}{-}0.085$ ($p{=}0.055$, NT-hg38) and $r{=}{-}0.102$ ($p{=}0.021$,
NT-1000G), explaining $R^2{=}$0.72\% and 1.05\%. Fisher's $z$-test confirms
the two within-group correlations are not significantly different
(NT-hg38: $p{=}0.48$; NT-1000G: $p{=}0.33$).

The standard deviation of bias scores (${\approx}2.23$--2.30 log-likelihood
units) vastly exceeds the between-group mean difference (0.43 units). The model's
reference preference is overwhelmingly determined by local sequence context, with
population allele frequency contributing less than 1\% of the variance.

\subsection{Local Variant Density Does Not Explain the Gap}

European-enriched variants reside in significantly denser genomic regions (mean
15.87 gnomAD neighbours within $\pm$500 bp) than African-enriched variants
(mean 7.93; $p{=}4.0{\times}10^{-26}$). Within the African-enriched group,
density correlates positively but marginally with bias ($r{=}{+}0.022$,
$p{=}0.054$), in the direction opposite to a context-complexity-driven
explanation. After matching African-enriched variants to the density range of
European-enriched variants (IQR 4--23 neighbours), the AFR bias mean changes
from $+$0.508 to $+$0.501; the gap with the EUR group ($+$0.941) is statistically
unchanged ($p{=}4.2{\times}10^{-5}$). Local genomic variant density is not a confounder.

\subsection{Downstream Scoring Asymmetry}

The scoring gap has a concrete directional consequence: 63.8\% of
African-enriched variants score below the European-enriched mean.
Negative bias (model preferring ALT over REF) affects 41.6\% of
African-enriched variants vs.\ 35.2\% of European-enriched variants
(odds ratio 1.31; $\chi^2{=}7.94$, $p{=}0.005$). This asymmetry runs in the
direction of \emph{reduced signal certainty} for African-enriched variants:
the model is less confident that the reference allele is preferred, producing
lower-magnitude scores that could translate to higher uncertainty or VUS
classifications in downstream pipelines that treat reference-allele preference
as a benignity signal. We measure a scoring difference only; demonstrating
propagation into clinical pathogenicity labels requires a separate experiment.

% ── V. Discussion ─────────────────────────────────────────────────────────────
\section{Discussion}

\subsection{The Asymmetry Is Resistant to Training-Set Diversification}

The 1.2\% change in the EUR--AFR gap when moving from NT-hg38 to NT-1000G is
the central finding of this paper. It indicates that the ancestry-associated
scoring asymmetry is resistant to training-set diversification under the
NT-1000G training formulation, and is consistent with a reference-centered
component in model scoring. One structural candidate: all training genomes in
NT-1000G---regardless of their ancestry---are expressed as sequences relative
to GRCh38. The model's local context at any scored position is therefore always
the linear reference haplotype, irrespective of how many diverse genomes were
seen during pretraining. Architecture, tokenization strategy, and haplotype
sampling choices may also contribute; our experiment does not isolate a single
cause but does show that training-genome diversity alone, in this formulation,
does not resolve the asymmetry.

This contrasts with protein language models, where diverse training sequences
can reduce ancestry bias \cite{pathak_2024} because protein LMs operate on
sequence families without a fixed linear reference coordinate system.

\subsection{Why Population Allele Frequency Is Not Encoded}

Both models show $R^2{<}1\%$ for within-group allele-frequency calibration.
The NT architecture has no explicit mechanism for encoding population-level
allele frequency; it learns sequence probability by aggregating over all training
genomic positions. The net effect is near-total insensitivity to population
frequency: a variant at $\text{AF}_\text{AFR}{=}30\%$ and one at 10\% are
nearly indistinguishable in their bias scores ($r{=}-0.058$). This is not a
failure of training diversity but a fundamental architectural limitation: without
explicit population labels or population-conditional training, MLM-based DNA LMs
cannot learn population allele frequencies.

\subsection{Direction of the Gap and Mechanism}

European-enriched variants receive \emph{higher} positive bias despite higher
global MAF---ruling out a frequency-calibration explanation. We propose that this
reflects differential haplotype context familiarity. European-enriched variants
reside in haplotype contexts that closely resemble the GRCh38 reference sequence;
the model has a strong, confident prior at these sites. African-enriched variants
sit in haplotype contexts that have diverged further from the European reference;
the model is less certain, manifesting as lower positive bias (closer to zero).
This hypothesis is consistent with the density analysis: AFR variants are in
sparser genomic regions (${\sim}2\times$ fewer common neighbours), yet density
itself does not account for the gap, suggesting the mechanism operates at the
haplotype sequence level rather than at variant density.

\subsection{Implications for Genomic AI and Health Equity}

Three implications follow. First, the observed scoring asymmetry warrants caution
when using NT-based bias scores as ancestry-neutral variant rankers. We measure a
\emph{scoring difference}, not a demonstrated clinical outcome; whether this
asymmetry propagates into downstream pathogenicity classifiers or VUS calls
requires a separate experiment directly linking bias scores to clinical labels
stratified by patient ancestry. Second, our result suggests that increasing
training-genome diversity under the current NT training formulation is
insufficient to close the scoring gap, implying that architectural changes or
reference-agnostic training are needed. Third, pangenome-based representations
(T2T-CHM13 \cite{t2t}, HPRC \cite{pangenome}) offer a promising direction by
removing the single-linear-reference anchor, but this remains to be tested
empirically in the foundation model setting.

\subsection{Limitations}

This study uses chromosome 22 only. Only NT models are evaluated; whether other
DNA LMs (DNABERT-2, HyenaDNA, Evo) show analogous patterns requires separate
investigation. We do not directly measure downstream clinical outcomes;
demonstrating that NT bias scores affect pathogenicity classifier performance
on African vs.\ European patient cohorts is an important future direction.

As a sensitivity analysis, we scored an expanded EUR set ($N_\text{EUR}{=}6{,}830$;
threshold $\text{AF}_\text{EUR}{\geq}0.05$). The EUR$-$AFR gap is directionally
consistent (NT-hg38: $\Delta{=}{+}0.215$, $p{=}5.3{\times}10^{-9}$, $d{=}0.097$;
NT-1000G: $\Delta{=}{+}0.204$, $p{=}3.4{\times}10^{-8}$, $d{=}0.092$), though
the expanded EUR set has a lower within-population AF mean (0.080 vs.\ AFR 0.179),
so the two sets are no longer matched for frequency magnitude. For this reason,
the $N{=}509$ AF-matched analysis is the primary result. In the expanded set,
Fisher's $z$-test detects a significant difference in within-group AF-calibration
correlations between AFR and EUR (NT-1000G: $p{=}6.7{\times}10^{-5}$), a pattern
that is underpowered in the matched set and warrants follow-up with a properly
matched larger EUR cohort.

% ── VI. Conclusion ────────────────────────────────────────────────────────────
\section{Conclusion}

We show that Nucleotide Transformer exhibits a significant ancestry-associated
scoring asymmetry: European-enriched common variants receive greater reference
allele preference than African-enriched variants at matched within-population
allele frequency ($\Delta{=}{+}0.43$, $p{<}5{\times}10^{-5}$, replicated
across both NT models). The EUR--AFR gap changes by only 1.2\% when pretraining
expands from a single reference genome to 3{,}202 diverse genomes. Neither model
encodes population allele frequency ($R^2{<}1\%$), and local variant density
does not account for the gap. These results indicate that the observed asymmetry
is resistant to training-set diversification under the NT-1000G formulation and
is consistent with a persistent reference-centered component in DNA
foundation-model scoring. Whether downstream clinical applications inherit this
asymmetry, and whether pangenome-based training resolves it, remain important
open questions for equitable genomic AI.

% ── Acknowledgements ──────────────────────────────────────────────────────────
\section*{Acknowledgements}
The authors thank the gnomAD consortium for making population-stratified allele
frequency data publicly available, and InstaDeepAI for releasing pre-trained
Nucleotide Transformer model weights under open licences. Code and scored data
are available at \url{https://github.com/thackshanaramana0-spec/nt_ancestry_bias}.

% ── References ────────────────────────────────────────────────────────────────
\bibliographystyle{IEEEtran}
\bibliography{references}

\end{document}
